\documentclass[11pt,a4paper]{article}
\usepackage[margin=25mm]{geometry}
\usepackage[T1]{fontenc}
\usepackage{lmodern}
\usepackage{amsmath,amssymb}
\usepackage{booktabs,array,makecell}
\usepackage[table]{xcolor}
\usepackage{float}
\usepackage{graphicx}
\usepackage[font=small]{caption}
\usepackage[hidelinks]{hyperref}
\usepackage{authblk}
\usepackage{listings}

\definecolor{fastc}{RGB}{0,128,0}
\definecolor{slowc}{RGB}{190,0,0}
\newcommand{\tfast}[1]{\textcolor{fastc}{#1}}
\newcommand{\tslow}[1]{\textcolor{slowc}{#1}}
\lstset{basicstyle=\ttfamily\footnotesize,columns=fullflexible,keepspaces=true,frame=single,framerule=0.3pt,
        rulecolor=\color{gray},language=C++,commentstyle=\color{gray!70!black},showstringspaces=false}

\title{TIGRE 3.3.1 and ASTRA 2.5:\\a speed comparison of the projection operators on one and several GPUs}
\author{Ander Biguri}
\affil{Department of Applied Mathematics and Theoretical Physics (DAMTP), University of Cambridge}
\date{November 2026\\[4pt]\small TIGRE technical note. Accompanies TIGRE 3.3.1 (CERN/TIGRE pull requests \#903--\#907 and \#911--\#917).}

\usepackage{adjustbox}
\begin{document}
\maketitle

\begin{abstract}
We compare the four projection operators of TIGRE 3.3.1 and ASTRA 2.5 --- forward projection (TIGRE's Siddon
and interpolated projectors against ASTRA's) and backprojection (TIGRE's matched and FDK-weighted
backprojectors against ASTRA's) --- on a desktop GPU (an RTX 4070, Windows) and on a server with four Quadro
RTX 6000 (Linux), for cone-beam, parallel-beam and fan-beam (2D) geometries, with host (NumPy) and GPU (PyTorch)
arrays, from $32^3$ to $1536^3$. For every case where ASTRA is faster we give the reason. On one GPU of the RTX
4070, TIGRE 3.3.1 is faster in 129 of 205 comparisons and ASTRA in 64, 51 of which are TIGRE's interpolated
projector, which was not optimised in this release; without it, TIGRE is faster in 127 of 151 and ASTRA in 13.
In 3D, the Siddon projector takes about half of ASTRA's time and the backprojectors 0.54--0.67 of it; in 2D the
backprojectors take a fifth, and the Siddon projector 1.0--1.3 times ASTRA's time, where ASTRA's interpolating
projector does less work per ray. On several GPUs, TIGRE gains less than ASTRA from extra GPUs for problems that
fit in one: its multi-GPU calls copy the whole image to every GPU (measured on one four-GPU server; other
machines may differ). For the large problems TIGRE's multi-GPU
design is for, its forward projection of a $1536^3$ image with 1536 angles takes 20.5\,s on one GPU and 10.4\,s
on four, against ASTRA's 83.8 and 24.2\,s.
\end{abstract}

\section{Introduction}
TIGRE~\cite{tigre2016,tigre2020multigpu,tigre2025v3} and ASTRA~\cite{astra2016} are the two open GPU toolboxes
most used for tomographic reconstruction from Python. Their operators are often compared, and the answer
depends on the problem: on the size of the image, on the geometry, on whether the data live on the host or on
the GPU, and on how many GPUs share the work. A comparison made for the release of TIGRE 3.3.0 measured large
cone-beam problems; this note covers the whole range, and it was made with a question in mind: where TIGRE is
slower, why?

Answering it led to the changes in TIGRE 3.3.1 (Section~\ref{sec:changes}). Most of them concern what a call does
besides its kernels: allocating and freeing memory, waiting for the GPU, page-locking host memory, copying the
geometry from Python. Those costs do not grow with the problem, so they decide the speed of small problems, of
2D (fan-beam) problems, of iterative algorithms, which call the operators hundreds of times, and of the
operators inside a neural network. A 2D fan-beam projection of a $256^2$ image with 360 angles took about 3\,ms in
TIGRE 3.3.0, of which 0.2\,ms were kernels.

Section~\ref{sec:setup} describes the machines and the method, Section~\ref{sec:results} the results, and
Section~\ref{sec:why} explains each case where ASTRA is faster.

\section{What changed in TIGRE 3.3.1}\label{sec:changes}
\begin{table}[H]
\centering\small
\caption{Changes in TIGRE 3.3.1 that affect speed. Results are bit-identical to 3.3.0's except where a
sum is ordered differently (parallel-beam backprojection, $2\times10^{-6}$ relative; fan-beam FDK filtering on
some CPUs, $3\times10^{-7}$).}
\label{tab:changes}
\begin{tabular}{>{\raggedright}p{0.17\textwidth}>{\raggedright\arraybackslash}p{0.75\textwidth}}
\toprule
pull request & change \\
\midrule
\#904 & Launch shapes for thin problems. Forward projectors: thread blocks along the long side of a detector of
fewer than 9 rows (a one-row fan-beam detector left 8 of 9 threads idle), and several blocks of angles per
launch. Cone backprojectors: the column of voxels per thread matched to images of fewer than 8 slices.\\
\#905, \#917 & One rule for page-locking the caller's arrays: only when the image is split or more than two
GPUs are used. On one GPU, and on two, locking an array cost more than it saved.\\
\#907 & The fixed cost of a call: the geometry copied from Python through typed memoryviews; the projection
buffers of the memory check used instead of allocated again; buffers cleared and textures filled on the
kernels' stream instead of waiting for the GPU; GPU arrays written to directly by the forward projectors; all
projections in one texture on thin detectors; the per-kernel constants in ordinary host memory.\\
\#911 & Cone backprojectors: no wait for the GPU between the kernels of a chunk of projections.\\
\#912 & Backprojectors read the projections from a layered 2D texture (bilinear fetch) instead of a 3D
texture (trilinear fetch at an exact projection index: the same value at half the texture rate). The
parallel-beam kernel gets the per-projection set-up folded on the host, as the cone kernels already had.\\
\#913, \#916 & Siddon forward projectors: several steps of the ray walk per pass, with selects instead of
branches, so that several texture fetches are in flight at once (the same sums): four for the usual detector,
fewer when the rays are far apart in a large image, where more fetches in flight evict each other from the
cache (\#916); and the constants of the next kernel launch computed on the host while the current kernels run,
instead of after waiting for them.\\
\#914 & FDK: detectors with few rows (fan beam), or a multiple of 16 rows that fits a block twice, filtered
several projections per FFT call; the cosine weights applied while the filtering copies each block into its FFT
buffer, instead of in a pass of their own.\\
\#915 & SIRT, SART, OS-SART, CGLS: vector updates in the arrays the operators return, without full-size
temporaries. The same arithmetic.\\
\#903, \#906 & Fixes: a crash in \texttt{Atb} at $1024^3$ with 32-bit sizes; page-locking that fails falls back to
pageable memory; a crash of the backprojection on 4+ GPUs when a GPU gets no slices.\\
\bottomrule
\end{tabular}
\end{table}
FDK and the iterative algorithms (\#914, \#915) are not compared in this note, which covers the operators only.

\section{Setup}\label{sec:setup}
\paragraph{Machines.}
Windows 11 (build 26300), NVIDIA GeForce RTX 4070 (12\,GB, compute capability 8.9), 32\,GB RAM, Python 3.11,
ASTRA 2.5.0. Linux server (Ubuntu 22.04, kernel 5.15): 4 NVIDIA Quadro RTX 6000 (24\,GB each, compute capability
7.5), 126\,GB RAM, CUDA 12.5, Python 3.12, PyTorch 2.6, ASTRA 2.5.0, TIGRE built for sm\_75. Both toolboxes are used through their Python interfaces:
TIGRE's \texttt{Ax} and \texttt{Atb}, and ASTRA's \texttt{direct\_FP3D} and \texttt{direct\_BP3D} from
\texttt{astra.experimental} (in 2D, \texttt{direct\_FP2D} and \texttt{direct\_BP2D}), which take host arrays or GPU tensors without copies on the Python side. Only the operators are compared:
reconstruction algorithms (FDK, SIRT, CGLS) depend on more than the operators (filtering, vector arithmetic,
where the data live) and are left out.

\paragraph{Same problems.}
The toolboxes use different conventions for angles, detector orientation and offsets. For each beam type the
suite searches the ASTRA convention that reproduces TIGRE's projections on a small problem and uses it
throughout (Table~\ref{tab:calib}), so that both compute the same projections of the same image: their results
agree to $5\times10^{-4}$ (relative $\ell_2$; TIGRE's Siddon projector integrates exactly, ASTRA's interpolates).

\paragraph{Timing.}
Each measurement is the median of several calls, repeated until it lasts at least 0.3\,s; the two toolboxes
alternate in rounds (three, the order changing every round), and the result is the ratio
$t_\text{TIGRE}/t_\text{ASTRA}$ with a 95\% confidence interval over the rounds. Below 1, TIGRE is faster
(\tfast{green}); above, ASTRA (\tslow{red}); $^\dagger$ marks an interval that contains 1. Host arrays include
the copies to and from the GPU; GPU arrays (PyTorch tensors) do not, for either toolbox. Changes to TIGRE were
measured the same way, alternating the two builds and using ASTRA, run in between, as a control for the
machine's drift.

\IfFileExists{tables/win_calib.tex}{\begin{table}[H]
\centering
\caption{Windows, RTX 4070. Calibration: the ASTRA convention that matches TIGRE, found on a small problem for each beam type, and how well the projections agree under it (relative $L_2$ difference after a least-squares scale). ``Next best'': the best convention that gives measurably different projections.}
\label{tab:calib}
\small
\begin{adjustbox}{max width=\linewidth}
\begin{tabular}{llrrr}
\toprule
beam & ASTRA convention & agree & scale & next best\\
\midrule
cone & $\theta_A = +1\,\theta_T +90^\circ$ & 5.4e-04 & 0.9998 & 1.6e-01\\
parallel & $\theta_A = +1\,\theta_T +270^\circ$, $u$ flipped & 5.1e-04 & 0.9998 & 4.9e-01\\
fan & $\theta_A = -1\,\theta_T +90^\circ$, $u$ flipped & 5.4e-04 & 0.9999 & 1.2e-01\\
offset & detector shift $+1\,u, +1\,v$ & 5.2e-04 & 0.9998 & 3.9e-01\\
\bottomrule
\end{tabular}
\end{adjustbox}
\end{table}
}{}

\section{Results}\label{sec:results}
\subsection{One GPU}
Tables~\ref{tab:summary} (RTX 4070) and~\ref{tab:lin-summary} (Quadro RTX 6000) summarise the comparisons on one
GPU; the tables that follow give every case on the RTX 4070, and Appendix~\ref{app:linux} those on the Quadro.
\IfFileExists{tables/win_summary.tex}{\begin{table}[H]
\centering\small
\caption{Summary, one GPU (RTX 4070, Windows): median ratio $t_\text{TIGRE}/t_\text{ASTRA}$ over the cases of
each class (below 1: TIGRE faster), its range, and the median speed-up of TIGRE 3.3.1 over 3.3.0 on the same
cases (``--'': not measured with 3.3.0).}
\label{tab:summary}
\begin{adjustbox}{max width=\linewidth}
\begin{tabular}{lllrrrr}
\toprule
operation & problems & arrays & cases & median & range & vs 3.3.0 \\
\midrule
fp Siddon & small 3D ($\le128^3$) & NumPy & 4 & \tfast{0.63} & 0.50--0.69 & 1.61$\times$ \\
fp Siddon & small 3D ($\le128^3$) & tensors & 4 & \tfast{0.65} & 0.53--0.92 & 2.30$\times$ \\
fp Siddon & cone $256^3$--$1024^3$ & NumPy & 15 & \tfast{0.46} & 0.29--0.81 & 1.10$\times$ \\
fp Siddon & cone $256^3$--$1024^3$ & tensors & 15 & \tfast{0.49} & 0.06--0.95 & 1.21$\times$ \\
fp Siddon & parallel & NumPy & 3 & \tfast{0.59} & 0.55--0.63 & -- \\
fp Siddon & parallel & tensors & 3 & \tfast{0.75} & 0.10--0.86 & -- \\
fp Siddon & fan (2D) & NumPy & 5 & \tslow{1.17} & 1.11--1.22 & 4.87$\times$ \\
fp Siddon & fan (2D) & tensors & 5 & \tslow{1.17} & 1.02--1.32 & 6.57$\times$ \\
\addlinespace[2pt]
fp interpolated & small 3D ($\le128^3$) & NumPy & 4 & \tslow{1.18} & 0.84--2.38 & 1.25$\times$ \\
fp interpolated & small 3D ($\le128^3$) & tensors & 4 & \tslow{1.68} & 1.06--4.68 & 1.43$\times$ \\
fp interpolated & cone $256^3$--$1024^3$ & NumPy & 15 & \tslow{2.64} & 1.85--3.09 & 0.97$\times$ \\
fp interpolated & cone $256^3$--$1024^3$ & tensors & 15 & \tslow{3.31} & 0.83--4.56 & 0.98$\times$ \\
fp interpolated & parallel & NumPy & 3 & \tslow{3.04} & 2.89--3.10 & -- \\
fp interpolated & parallel & tensors & 3 & \tslow{4.44} & 1.08--4.89 & -- \\
fp interpolated & fan (2D) & NumPy & 5 & \tslow{5.76} & 3.20--6.79 & 2.32$\times$ \\
fp interpolated & fan (2D) & tensors & 5 & \tslow{7.23} & 3.05--8.17 & 2.46$\times$ \\
\addlinespace[2pt]
bp matched & small 3D ($\le128^3$) & NumPy & 4 & \tfast{0.85} & 0.55--0.92 & 1.75$\times$ \\
bp matched & small 3D ($\le128^3$) & tensors & 4 & \tslow{1.23} & 0.73--1.71 & 1.87$\times$ \\
bp matched & cone $256^3$--$1024^3$ & NumPy & 15 & \tfast{0.65} & 0.50--0.90 & 1.37$\times$ \\
bp matched & cone $256^3$--$1024^3$ & tensors & 13 & \tfast{0.67} & 0.54--0.77 & 1.47$\times$ \\
bp matched & fan (2D) & NumPy & 5 & \tfast{0.21} & 0.19--0.43 & 5.88$\times$ \\
bp matched & fan (2D) & tensors & 5 & \tfast{0.18} & 0.15--0.36 & 7.22$\times$ \\
\addlinespace[2pt]
bp FDK weights & small 3D ($\le128^3$) & NumPy & 4 & \tfast{0.87} & 0.56--1.02 & 1.71$\times$ \\
bp FDK weights & small 3D ($\le128^3$) & tensors & 4 & \tslow{1.29} & 0.69--1.66 & 1.84$\times$ \\
bp FDK weights & cone $256^3$--$1024^3$ & NumPy & 15 & \tfast{0.59} & 0.45--0.88 & 1.52$\times$ \\
bp FDK weights & cone $256^3$--$1024^3$ & tensors & 13 & \tfast{0.54} & 0.51--0.65 & 1.64$\times$ \\
bp FDK weights & fan (2D) & NumPy & 5 & \tfast{0.19} & 0.18--0.42 & 6.16$\times$ \\
bp FDK weights & fan (2D) & tensors & 5 & \tfast{0.16} & 0.13--0.34 & 7.80$\times$ \\
\addlinespace[2pt]
bp & parallel & NumPy & 3 & \tfast{0.55} & 0.49--0.61 & -- \\
bp & parallel & tensors & 2 & \tfast{0.53} & 0.52--0.53 & -- \\
\bottomrule
\end{tabular}
\end{adjustbox}
\end{table}
}{}
\IfFileExists{tables/lin_summary.tex}{\begin{table}[H]
\centering\small
\caption{Summary, one GPU (Quadro RTX 6000, Linux): median ratio $t_\text{TIGRE}/t_\text{ASTRA}$ over the cases of
each class (below 1: TIGRE faster), its range, and the median speed-up of TIGRE 3.3.1 over 3.3.0 on the same
cases (``--'': not measured with 3.3.0).}
\label{tab:lin-summary}
\begin{adjustbox}{max width=\linewidth}
\begin{tabular}{lllrrrr}
\toprule
operation & problems & arrays & cases & median & range & vs 3.3.0 \\
\midrule
fp Siddon & small 3D ($\le128^3$) & NumPy & 4 & \tfast{0.61} & 0.54--0.69 & 1.75$\times$ \\
fp Siddon & small 3D ($\le128^3$) & tensors & 4 & \tfast{0.70} & 0.52--0.83 & 2.29$\times$ \\
fp Siddon & cone $256^3$--$1024^3$ & NumPy & 15 & \tfast{0.48} & 0.40--0.81 & 1.74$\times$ \\
fp Siddon & cone $256^3$--$1024^3$ & tensors & 15 & \tfast{0.49} & 0.42--0.90 & 2.17$\times$ \\
fp Siddon & parallel & NumPy & 3 & \tfast{0.66} & 0.60--0.94 & 1.81$\times$ \\
fp Siddon & parallel & tensors & 3 & \tfast{0.70} & 0.69--0.84 & 2.05$\times$ \\
fp Siddon & fan (2D) & NumPy & 5 & \tslow{1.67} & 1.47--2.25 & 8.41$\times$ \\
fp Siddon & fan (2D) & tensors & 5 & \tslow{1.59} & 1.49--2.20 & 10.01$\times$ \\
\addlinespace[2pt]
fp interpolated & small 3D ($\le128^3$) & NumPy & 4 & 1.03 & 0.80--2.06 & 1.36$\times$ \\
fp interpolated & small 3D ($\le128^3$) & tensors & 4 & \tslow{1.30} & 1.01--3.18 & 1.55$\times$ \\
fp interpolated & cone $256^3$--$1024^3$ & NumPy & 15 & \tslow{2.60} & 1.78--3.13 & 1.00$\times$ \\
fp interpolated & cone $256^3$--$1024^3$ & tensors & 15 & \tslow{3.34} & 2.78--3.79 & 1.01$\times$ \\
fp interpolated & parallel & NumPy & 3 & \tslow{2.96} & 2.66--3.27 & 1.01$\times$ \\
fp interpolated & parallel & tensors & 3 & \tslow{4.17} & 4.12--4.72 & 1.01$\times$ \\
fp interpolated & fan (2D) & NumPy & 5 & \tslow{6.04} & 4.18--7.85 & 2.50$\times$ \\
fp interpolated & fan (2D) & tensors & 5 & \tslow{7.41} & 4.31--8.30 & 2.54$\times$ \\
\addlinespace[2pt]
bp matched & small 3D ($\le128^3$) & NumPy & 4 & \tfast{0.88} & 0.72--0.90 & 1.70$\times$ \\
bp matched & small 3D ($\le128^3$) & tensors & 4 & \tslow{1.10} & 0.91--1.13 & 1.79$\times$ \\
bp matched & cone $256^3$--$1024^3$ & NumPy & 15 & \tslow{1.08} & 0.64--1.75 & 1.39$\times$ \\
bp matched & cone $256^3$--$1024^3$ & tensors & 15 & \tfast{0.80} & 0.76--0.92 & 1.51$\times$ \\
bp matched & fan (2D) & NumPy & 5 & \tfast{0.33} & 0.25--0.74 & 5.61$\times$ \\
bp matched & fan (2D) & tensors & 5 & \tfast{0.22} & 0.19--0.63 & 7.69$\times$ \\
\addlinespace[2pt]
bp FDK weights & small 3D ($\le128^3$) & NumPy & 4 & \tfast{0.87} & 0.71--0.90 & 1.69$\times$ \\
bp FDK weights & small 3D ($\le128^3$) & tensors & 4 & \tslow{1.04} & 0.79--1.12 & 1.83$\times$ \\
bp FDK weights & cone $256^3$--$1024^3$ & NumPy & 15 & 0.98 & 0.67--1.73 & 1.40$\times$ \\
bp FDK weights & cone $256^3$--$1024^3$ & tensors & 15 & \tfast{0.60} & 0.59--0.77 & 1.60$\times$ \\
bp FDK weights & fan (2D) & NumPy & 5 & \tfast{0.31} & 0.24--0.73 & 5.13$\times$ \\
bp FDK weights & fan (2D) & tensors & 5 & \tfast{0.20} & 0.17--0.62 & 6.77$\times$ \\
\addlinespace[2pt]
bp & parallel & NumPy & 3 & 0.98 & 0.87--0.98 & 1.66$\times$ \\
bp & parallel & tensors & 3 & \tfast{0.54} & 0.54--0.57 & 2.46$\times$ \\
\bottomrule
\end{tabular}
\end{adjustbox}
\end{table}
}{}
\IfFileExists{tables/win_size1.tex}{\begin{table}[H]
\centering
\caption{Windows, RTX 4070. Cone beam, $n^3$ image, $n^2$ detector covering the image (409.6\,mm), 360 angles, 1 GPU.}
\label{tab:size1}
\small\setlength{\tabcolsep}{4pt}
\begin{adjustbox}{max width=\linewidth}
\begin{tabular}{llrrrrrr}
\toprule
& & \multicolumn{3}{c}{NumPy arrays, 1 GPU} & \multicolumn{3}{c}{PyTorch GPU tensors, 1 GPU}\\
\cmidrule(lr){3-5}\cmidrule(lr){6-8}
problem & operation & TIGRE ms & ASTRA ms & ratio & TIGRE ms & ASTRA ms & ratio\\
\midrule
128$^3$ & fp, Siddon & 7.90 & 11.5 & \tfast{$0.68\pm0.09$} & 5.35 & 5.81 & \tfast{$0.93\pm0.10$}$^\dagger$\\
 & fp, interpolated & 27.3 & 11.5 & \tslow{$2.36\pm0.06$} & 27.2 & 5.81 & \tslow{$4.65\pm0.20$}\\
 & bp, matched & 5.21 & 9.44 & \tfast{$0.56\pm0.06$} & 3.52 & 4.80 & \tfast{$0.76\pm0.13$}\\
 & bp, FDK weights & 5.25 & 9.44 & \tfast{$0.57\pm0.11$} & 3.31 & 4.80 & \tfast{$0.70\pm0.07$}\\
\addlinespace
256$^3$ & fp, Siddon & 32.1 & 63.9 & \tfast{$0.50\pm0.05$} & 27.1 & 38.7 & \tfast{$0.70\pm0.05$}\\
 & fp, interpolated & 181 & 63.9 & \tslow{$2.86\pm0.23$} & 177 & 38.7 & \tslow{$4.55\pm0.07$}\\
 & bp, matched & 32.5 & 54.7 & \tfast{$0.58\pm0.05$} & 21.3 & 33.0 & \tfast{$0.65\pm0.03$}\\
 & bp, FDK weights & 29.1 & 54.7 & \tfast{$0.53\pm0.01$} & 18.4 & 33.0 & \tfast{$0.56\pm0.02$}\\
\addlinespace
512$^3$ & fp, Siddon & 241 & 526 & \tfast{$0.46\pm0.01$} & 202 & 407 & \tfast{$0.50\pm0.01$}\\
 & fp, interpolated & 1384 & 526 & \tslow{$2.63\pm0.08$} & 1366 & 407 & \tslow{$3.35\pm0.06$}\\
 & bp, matched & 242 & 375 & \tfast{$0.65\pm0.01$} & 163 & 253 & \tfast{$0.64\pm0.01$}\\
 & bp, FDK weights & 219 & 375 & \tfast{$0.59\pm0.03$} & 136 & 253 & \tfast{$0.54\pm0.02$}\\
\addlinespace
1024$^3$ & fp, Siddon & 2137 & 4549 & \tfast{$0.48\pm0.04$} & 1697 & 24480 & \tfast{$0.07\pm0.01$}\\
 & fp, interpolated & 11280 & 4549 & \tslow{$2.48\pm0.02$} & 20439 & 24480 & \tfast{$0.86\pm0.10$}\\
 & bp, matched & 2076 & 2964 & \tfast{$0.73\pm0.18$} & --- & --- & ---\\
 & bp, FDK weights & 1825 & 2964 & \tfast{$0.64\pm0.07$} & --- & --- & ---\\
\bottomrule
\end{tabular}
\end{adjustbox}
\end{table}
}{}
\IfFileExists{tables/win_legacy.tex}{\begin{table}[H]
\centering
\caption{Windows, RTX 4070. The 3.3.0 note's detector (0.8\,mm pixels whatever $n$), 1 GPU. Compare with Table~\ref{tab:size1} at the same $n$.}
\label{tab:legacy}
\small\setlength{\tabcolsep}{4pt}
\begin{adjustbox}{max width=\linewidth}
\begin{tabular}{llrrrrrr}
\toprule
& & \multicolumn{3}{c}{NumPy arrays, 1 GPU} & \multicolumn{3}{c}{PyTorch GPU tensors, 1 GPU}\\
\cmidrule(lr){3-5}\cmidrule(lr){6-8}
problem & operation & TIGRE ms & ASTRA ms & ratio & TIGRE ms & ASTRA ms & ratio\\
\midrule
256$^3$, 0.8 mm & fp, Siddon & 31.5 & 56.8 & \tfast{$0.55\pm0.05$} & 27.4 & 37.3 & \tfast{$0.73\pm0.04$}\\
 & fp, interpolated & 175 & 56.8 & \tslow{$3.08\pm0.17$} & 170 & 37.3 & \tslow{$4.57\pm0.14$}\\
 & bp, matched & 31.9 & 54.9 & \tfast{$0.58\pm0.04$} & 21.3 & 32.6 & \tfast{$0.64\pm0.03$}\\
 & bp, FDK weights & 28.9 & 54.9 & \tfast{$0.53\pm0.08$} & 18.0 & 32.6 & \tfast{$0.56\pm0.01$}\\
\addlinespace
1024$^3$, 0.8 mm & fp, Siddon & 1352 & 4276 & \tfast{$0.32\pm0.02$} & 1013 & 17949 & \tfast{$0.06\pm0.00$}\\
 & fp, interpolated & 12300 & 4276 & \tslow{$2.87\pm0.07$} & 22654 & 17949 & \tslow{$1.27\pm0.05$}\\
 & bp, matched & 1959 & 2655 & \tfast{$0.74\pm0.03$} & --- & --- & ---\\
 & bp, FDK weights & 1739 & 2655 & \tfast{$0.67\pm0.05$} & --- & --- & ---\\
\bottomrule
\end{tabular}
\end{adjustbox}
\end{table}
}{}
\IfFileExists{tables/win_angles.tex}{\begin{table}[H]
\centering
\caption{Windows, RTX 4070. Cone beam $512^3$, varying the number of angles, 1 GPU.}
\label{tab:angles}
\small\setlength{\tabcolsep}{4pt}
\begin{adjustbox}{max width=\linewidth}
\begin{tabular}{llrrrrrr}
\toprule
& & \multicolumn{3}{c}{NumPy arrays, 1 GPU} & \multicolumn{3}{c}{PyTorch GPU tensors, 1 GPU}\\
\cmidrule(lr){3-5}\cmidrule(lr){6-8}
problem & operation & TIGRE ms & ASTRA ms & ratio & TIGRE ms & ASTRA ms & ratio\\
\midrule
45 angles & fp, Siddon & 82.5 & 119 & \tfast{$0.69\pm0.10$} & 43.2 & 55.6 & \tfast{$0.79\pm0.07$}\\
 & fp, interpolated & 219 & 119 & \tslow{$1.83\pm0.25$} & 181 & 55.6 & \tslow{$3.27\pm0.04$}\\
 & bp, matched & 111 & 123 & \tfast{$0.91\pm0.06$} & 33.0 & 42.8 & \tfast{$0.77\pm0.02$}\\
 & bp, FDK weights & 108 & 123 & \tfast{$0.91\pm0.20$}$^\dagger$ & 27.9 & 42.8 & \tfast{$0.66\pm0.01$}\\
\addlinespace
90 angles & fp, Siddon & 105 & 162 & \tfast{$0.67\pm0.22$} & 67.5 & 93.5 & \tfast{$0.72\pm0.03$}\\
 & fp, interpolated & 380 & 162 & \tslow{$2.30\pm0.22$} & 344 & 93.5 & \tslow{$3.69\pm0.07$}\\
 & bp, matched & 128 & 154 & \tfast{$0.87\pm0.29$}$^\dagger$ & 50.6 & 70.3 & \tfast{$0.72\pm0.01$}\\
 & bp, FDK weights & 120 & 154 & \tfast{$0.78\pm0.12$} & 42.8 & 70.3 & \tfast{$0.60\pm0.01$}\\
\addlinespace
180 angles & fp, Siddon & 147 & 284 & \tfast{$0.52\pm0.04$} & 111 & 198 & \tfast{$0.56\pm0.02$}\\
 & fp, interpolated & 708 & 284 & \tslow{$2.49\pm0.01$} & 672 & 198 & \tslow{$3.39\pm0.04$}\\
 & bp, matched & 181 & 227 & \tfast{$0.77\pm0.09$} & 88.6 & 129 & \tfast{$0.69\pm0.02$}\\
 & bp, FDK weights & 150 & 227 & \tfast{$0.65\pm0.05$} & 72.9 & 129 & \tfast{$0.57\pm0.00$}\\
\addlinespace
360 angles & fp, Siddon & 232 & 519 & \tfast{$0.45\pm0.01$} & 198 & 400 & \tfast{$0.50\pm0.02$}\\
 & fp, interpolated & 1372 & 519 & \tslow{$2.64\pm0.11$} & 1324 & 400 & \tslow{$3.31\pm0.01$}\\
 & bp, matched & 243 & 364 & \tfast{$0.67\pm0.10$} & 165 & 247 & \tfast{$0.66\pm0.02$}\\
 & bp, FDK weights & 217 & 364 & \tfast{$0.59\pm0.02$} & 134 & 247 & \tfast{$0.54\pm0.02$}\\
\addlinespace
720 angles & fp, Siddon & 401 & 944 & \tfast{$0.43\pm0.02$} & 374 & 810 & \tfast{$0.47\pm0.01$}\\
 & fp, interpolated & 2655 & 944 & \tslow{$2.81\pm0.06$} & 2635 & 810 & \tslow{$3.25\pm0.01$}\\
 & bp, matched & 411 & 632 & \tfast{$0.64\pm0.04$} & 321 & 481 & \tfast{$0.66\pm0.02$}\\
 & bp, FDK weights & 333 & 632 & \tfast{$0.53\pm0.03$} & 255 & 481 & \tfast{$0.53\pm0.01$}\\
\addlinespace
1440 angles & fp, Siddon & 755 & 1861 & \tfast{$0.40\pm0.01$} & 731 & 1531 & \tfast{$0.48\pm0.01$}\\
 & fp, interpolated & 5290 & 1861 & \tslow{$2.83\pm0.04$} & 5280 & 1531 & \tslow{$3.44\pm0.04$}\\
 & bp, matched & 721 & 1223 & \tfast{$0.59\pm0.01$} & 644 & 950 & \tfast{$0.68\pm0.01$}\\
 & bp, FDK weights & 599 & 1223 & \tfast{$0.48\pm0.02$} & 501 & 950 & \tfast{$0.53\pm0.01$}\\
\bottomrule
\end{tabular}
\end{adjustbox}
\end{table}
}{}
\IfFileExists{tables/win_shape.tex}{\begin{table}[H]
\centering
\caption{Windows, RTX 4070. Non-cubic problems, 1 GPU: clinical CBCT ($256\times512\times512$ image, $512\times1024$ detector), fine detector ($256^3$ image, $1024^2$ detector), coarse detector ($512^3$ image, $256^2$ detector), 360 angles.}
\label{tab:shape}
\small\setlength{\tabcolsep}{4pt}
\begin{adjustbox}{max width=\linewidth}
\begin{tabular}{llrrrrrr}
\toprule
& & \multicolumn{3}{c}{NumPy arrays, 1 GPU} & \multicolumn{3}{c}{PyTorch GPU tensors, 1 GPU}\\
\cmidrule(lr){3-5}\cmidrule(lr){6-8}
problem & operation & TIGRE ms & ASTRA ms & ratio & TIGRE ms & ASTRA ms & ratio\\
\midrule
clinical CBCT & fp, Siddon & 298 & 896 & \tfast{$0.33\pm0.01$} & 288 & 774 & \tfast{$0.37\pm0.01$}\\
 & fp, interpolated & 2532 & 896 & \tslow{$2.81\pm0.09$} & 2521 & 774 & \tslow{$3.26\pm0.02$}\\
 & bp, matched & 139 & 276 & \tfast{$0.52\pm0.12$} & 93.9 & 146 & \tfast{$0.64\pm0.00$}\\
 & bp, FDK weights & 126 & 276 & \tfast{$0.45\pm0.02$} & 78.0 & 146 & \tfast{$0.54\pm0.01$}\\
\addlinespace
fine detector & fp, Siddon & 302 & 1035 & \tfast{$0.30\pm0.06$} & 306 & 773 & \tfast{$0.40\pm0.02$}\\
 & fp, interpolated & 2408 & 1035 & \tslow{$2.32\pm0.22$} & 2419 & 773 & \tslow{$3.13\pm0.01$}\\
 & bp, matched & 138 & 257 & \tfast{$0.54\pm0.02$} & 42.6 & 79.4 & \tfast{$0.54\pm0.01$}\\
 & bp, FDK weights & 144 & 257 & \tfast{$0.56\pm0.23$} & 40.4 & 79.4 & \tfast{$0.51\pm0.01$}\\
\addlinespace
coarse detector & fp, Siddon & 142 & 174 & \tfast{$0.82\pm0.05$} & 101 & 106 & \tfast{$0.95\pm0.01$}\\
 & fp, interpolated & 445 & 174 & \tslow{$2.54\pm0.10$} & 407 & 106 & \tslow{$3.82\pm0.05$}\\
 & bp, matched & 233 & 322 & \tfast{$0.72\pm0.07$} & 161 & 234 & \tfast{$0.69\pm0.05$}\\
 & bp, FDK weights & 222 & 322 & \tfast{$0.68\pm0.11$} & 132 & 234 & \tfast{$0.57\pm0.03$}\\
\bottomrule
\end{tabular}
\end{adjustbox}
\end{table}
}{}
\IfFileExists{tables/win_offset.tex}{\begin{table}[H]
\centering
\caption{Windows, RTX 4070. Cone beam $512^3$ with the detector shifted sideways by 102.4\,mm (a quarter of its width; ASTRA \texttt{cone\_vec}), 1 GPU.}
\label{tab:offset}
\small\setlength{\tabcolsep}{4pt}
\begin{adjustbox}{max width=\linewidth}
\begin{tabular}{llrrrrrr}
\toprule
& & \multicolumn{3}{c}{NumPy arrays, 1 GPU} & \multicolumn{3}{c}{PyTorch GPU tensors, 1 GPU}\\
\cmidrule(lr){3-5}\cmidrule(lr){6-8}
problem & operation & TIGRE ms & ASTRA ms & ratio & TIGRE ms & ASTRA ms & ratio\\
\midrule
u offset 102.4 mm & fp, Siddon & 203 & 507 & \tfast{$0.39\pm0.02$} & 167 & 398 & \tfast{$0.42\pm0.00$}\\
 & fp, interpolated & 1349 & 507 & \tslow{$2.65\pm0.06$} & 1318 & 398 & \tslow{$3.31\pm0.01$}\\
 & bp, matched & 243 & 381 & \tfast{$0.66\pm0.20$} & 162 & 246 & \tfast{$0.66\pm0.01$}\\
 & bp, FDK weights & 222 & 381 & \tfast{$0.58\pm0.06$} & 133 & 246 & \tfast{$0.54\pm0.00$}\\
\bottomrule
\end{tabular}
\end{adjustbox}
\end{table}
}{}
\IfFileExists{tables/win_parallel1.tex}{\begin{table}[H]
\centering
\caption{Windows, RTX 4070. Parallel beam (TIGRE parallel mode, ASTRA \texttt{parallel3d}), $n^3$, 360 angles, 1 GPU.}
\label{tab:parallel1}
\small\setlength{\tabcolsep}{4pt}
\begin{adjustbox}{max width=\linewidth}
\begin{tabular}{llrrrrrr}
\toprule
& & \multicolumn{3}{c}{NumPy arrays, 1 GPU} & \multicolumn{3}{c}{PyTorch GPU tensors, 1 GPU}\\
\cmidrule(lr){3-5}\cmidrule(lr){6-8}
problem & operation & TIGRE ms & ASTRA ms & ratio & TIGRE ms & ASTRA ms & ratio\\
\midrule
256$^3$ & fp, Siddon & 30.1 & 54.8 & \tfast{$0.55\pm0.02$} & 27.0 & 31.4 & \tfast{$0.85\pm0.03$}\\
 & fp, interpolated & 158 & 54.8 & \tslow{$2.90\pm0.06$} & 154 & 31.4 & \tslow{$4.87\pm0.06$}\\
 & bp & 27.7 & 56.4 & \tfast{$0.49\pm0.05$} & 16.9 & 32.2 & \tfast{$0.53\pm0.00$}\\
\addlinespace
512$^3$ & fp, Siddon & 230 & 388 & \tfast{$0.60\pm0.08$} & 196 & 263 & \tfast{$0.74\pm0.01$}\\
 & fp, interpolated & 1202 & 388 & \tslow{$3.12\pm0.11$} & 1166 & 263 & \tslow{$4.43\pm0.03$}\\
 & bp & 200 & 362 & \tfast{$0.57\pm0.12$} & 124 & 236 & \tfast{$0.53\pm0.01$}\\
\addlinespace
1024$^3$ & fp, Siddon & 1965 & 3120 & \tfast{$0.64\pm0.03$} & 1601 & 16762 & \tfast{$0.10\pm0.00$}\\
 & fp, interpolated & 9482 & 3120 & \tslow{$3.03\pm0.12$} & 18039 & 16762 & \tslow{$1.08\pm0.00$}\\
 & bp & 1583 & 2590 & \tfast{$0.64\pm0.12$} & --- & --- & ---\\
\bottomrule
\end{tabular}
\end{adjustbox}
\end{table}
}{}
\IfFileExists{tables/win_fan.tex}{\begin{table}[H]
\centering
\caption{Windows, RTX 4070. Fan beam (2D): TIGRE's one-slice cone geometry against ASTRA's \texttt{fanflat} (FP\_CUDA, BP\_CUDA, FBP\_CUDA), $n^2$ image, $n$ detector pixels, 1 GPU.}
\label{tab:fan}
\small\setlength{\tabcolsep}{3pt}
\begin{adjustbox}{max width=\linewidth}
\begin{tabular}{llrrrrrr}
\toprule
& & \multicolumn{3}{c}{NumPy arrays, 1 GPU} & \multicolumn{3}{c}{PyTorch GPU tensors, 1 GPU}\\
\cmidrule(lr){3-5}\cmidrule(lr){6-8}
problem & operation & TIGRE ms & ASTRA ms & ratio & TIGRE ms & ASTRA ms & ratio\\
\midrule
512$^2$ & fp, Siddon & 1.51 & 1.28 & \tslow{$1.17\pm0.02$} & 1.23 & 1.21 & \tslow{$1.09\pm0.16$}$^\dagger$\\
 & fp, interpolated & 4.11 & 1.28 & \tslow{$3.31\pm0.69$} & 3.68 & 1.21 & \tslow{$3.15\pm0.43$}\\
 & bp, matched & 1.56 & 3.62 & \tfast{$0.43\pm0.03$} & 1.22 & 3.41 & \tfast{$0.36\pm0.00$}\\
 & bp, FDK weights & 1.51 & 3.62 & \tfast{$0.44\pm0.09$} & 1.16 & 3.41 & \tfast{$0.34\pm0.00$}\\
\addlinespace
1024$^2$ & fp, Siddon & 2.99 & 2.70 & \tslow{$1.10\pm0.03$} & 2.15 & 2.09 & \tslow{$1.04\pm0.03$}\\
 & fp, interpolated & 11.3 & 2.70 & \tslow{$4.17\pm0.09$} & 10.5 & 2.09 & \tslow{$5.04\pm0.14$}\\
 & bp, matched & 3.25 & 12.1 & \tfast{$0.27\pm0.01$} & 2.42 & 11.5 & \tfast{$0.21\pm0.00$}\\
 & bp, FDK weights & 3.11 & 12.1 & \tfast{$0.26\pm0.02$} & 2.25 & 11.5 & \tfast{$0.20\pm0.00$}\\
\addlinespace
2048$^2$ & fp, Siddon & 7.69 & 6.90 & \tslow{$1.12\pm0.03$} & 6.21 & 5.30 & \tslow{$1.17\pm0.03$}\\
 & fp, interpolated & 39.7 & 6.90 & \tslow{$5.79\pm0.24$} & 38.3 & 5.30 & \tslow{$7.22\pm0.03$}\\
 & bp, matched & 9.64 & 46.6 & \tfast{$0.21\pm0.00$} & 7.20 & 44.5 & \tfast{$0.16\pm0.00$}\\
 & bp, FDK weights & 8.96 & 46.6 & \tfast{$0.19\pm0.00$} & 6.39 & 44.5 & \tfast{$0.14\pm0.00$}\\
\addlinespace
4096$^2$ & fp, Siddon & 29.7 & 24.8 & \tslow{$1.17\pm0.11$} & 24.2 & 18.8 & \tslow{$1.29\pm0.03$}\\
 & fp, interpolated & 158 & 24.8 & \tslow{$6.30\pm0.28$} & 154 & 18.8 & \tslow{$8.19\pm0.13$}\\
 & bp, matched & 36.7 & 185 & \tfast{$0.20\pm0.01$} & 26.7 & 175 & \tfast{$0.15\pm0.00$}\\
 & bp, FDK weights & 33.1 & 185 & \tfast{$0.18\pm0.01$} & 23.1 & 175 & \tfast{$0.13\pm0.00$}\\
\addlinespace
1024$^2$, 1440 angles & fp, Siddon & 7.13 & 5.83 & \tslow{$1.22\pm0.04$} & 6.67 & 5.05 & \tslow{$1.32\pm0.01$}\\
 & fp, interpolated & 39.6 & 5.83 & \tslow{$6.78\pm0.22$} & 39.3 & 5.05 & \tslow{$7.78\pm0.05$}\\
 & bp, matched & 8.79 & 45.1 & \tfast{$0.19\pm0.00$} & 7.94 & 44.5 & \tfast{$0.18\pm0.01$}\\
 & bp, FDK weights & 8.03 & 45.1 & \tfast{$0.18\pm0.00$} & 7.16 & 44.5 & \tfast{$0.16\pm0.01$}\\
\bottomrule
\end{tabular}
\end{adjustbox}
\end{table}
}{}
\IfFileExists{tables/win_latency.tex}{\begin{table}[H]
\centering
\caption{Windows, RTX 4070. Small problems (cone beam, $n^3$ image, $n^2$ detector, 100 angles), 1 GPU.}
\label{tab:latency}
\small\setlength{\tabcolsep}{4pt}
\begin{adjustbox}{max width=\linewidth}
\begin{tabular}{llrrrrrr}
\toprule
& & \multicolumn{3}{c}{NumPy arrays, 1 GPU} & \multicolumn{3}{c}{PyTorch GPU tensors, 1 GPU}\\
\cmidrule(lr){3-5}\cmidrule(lr){6-8}
problem & operation & TIGRE ms & ASTRA ms & ratio & TIGRE ms & ASTRA ms & ratio\\
\midrule
32$^3$ & fp, Siddon & 1.33 & 1.94 & \tfast{$0.68\pm0.05$} & 0.97 & 1.26 & \tfast{$0.79\pm0.21$}$^\dagger$\\
 & fp, interpolated & 1.64 & 1.94 & \tfast{$0.86\pm0.10$} & 1.33 & 1.26 & \tslow{$1.08\pm0.05$}\\
 & bp, matched & 0.98 & 1.07 & \tfast{$0.94\pm0.14$}$^\dagger$ & 0.94 & 0.55 & \tslow{$1.64\pm0.35$}\\
 & bp, FDK weights & 1.10 & 1.07 & \tslow{$1.00\pm0.21$}$^\dagger$ & 0.91 & 0.55 & \tslow{$1.58\pm0.31$}\\
\addlinespace
64$^3$ & fp, Siddon & 1.71 & 2.95 & \tfast{$0.57\pm0.09$} & 1.01 & 1.91 & \tfast{$0.53\pm0.01$}\\
 & fp, interpolated & 2.75 & 2.95 & \tfast{$0.92\pm0.10$}$^\dagger$ & 2.32 & 1.91 & \tslow{$1.26\pm0.23$}\\
 & bp, matched & 1.28 & 1.39 & \tfast{$0.97\pm0.22$}$^\dagger$ & 0.92 & 0.62 & \tslow{$1.51\pm0.02$}\\
 & bp, FDK weights & 1.33 & 1.39 & \tfast{$0.95\pm0.08$}$^\dagger$ & 1.00 & 0.62 & \tslow{$1.58\pm0.27$}\\
\addlinespace
128$^3$ & fp, Siddon & 2.94 & 5.86 & \tfast{$0.50\pm0.01$} & 1.88 & 3.56 & \tfast{$0.53\pm0.02$}\\
 & fp, interpolated & 8.36 & 5.86 & \tslow{$1.44\pm0.09$} & 7.64 & 3.56 & \tslow{$2.14\pm0.02$}\\
 & bp, matched & 3.19 & 4.03 & \tfast{$0.80\pm0.07$} & 1.63 & 1.68 & \tfast{$0.97\pm0.04$}$^\dagger$\\
 & bp, FDK weights & 3.15 & 4.03 & \tfast{$0.79\pm0.03$} & 1.61 & 1.68 & \tfast{$0.99\pm0.18$}$^\dagger$\\
\bottomrule
\end{tabular}
\end{adjustbox}
\end{table}
}{}

\paragraph{The fixed cost of a call.}
Table~\ref{tab:fixedcost} follows one fan-beam call through the releases. In 3.3.0 a forward projection of a
$256^2$ image with 360 angles took 7.9\,ms, of which 0.2\,ms were kernels; the rest was allocating and freeing
GPU and page-locked memory, waiting for the GPU, and converting the geometry in Python, all of which cost far
more under Windows' WDDM driver model than the kernels.

\begin{table}[H]
\centering\small
\caption{Wall time of one fan-beam call ($256^2$ image, 360 angles, PyTorch tensors, RTX 4070), from Nsight
Systems traces, in ms.}
\label{tab:fixedcost}
\begin{tabular}{lrr}
\toprule
version & Siddon fp & matched bp \\
\midrule
3.3.0 & 7.9 & -- \\
+ launch shapes (\#904) & 3.0 & 3.0 \\
+ fixed-cost changes (\#907) & 1.13 & 1.39 \\
+ no wait between kernels (\#911) & 1.13 & 1.29 \\
ASTRA 2.5 & 1.03 & 1.35 \\
\bottomrule
\end{tabular}
\end{table}

\subsection{Several GPUs}
Tables~\ref{tab:lin-size_multi} and~\ref{tab:lin-parallel_multi} give the times on 1, 2 and 4 GPUs of the Linux
server, with NumPy arrays. Every problem in these tables fits in one 24\,GB GPU. At $1024^3$, two GPUs are faster
than one for every TIGRE operator: the Siddon forward projection takes 1.79\,s on one GPU and 1.44\,s on two, faster
than ASTRA on two (2.17\,s), and the matched backprojection 3.03 and 2.36\,s. At $512^3$ two GPUs are faster for all
but the Siddon projector, and at $256^3$ only for the interpolated projector. On four GPUs TIGRE is slower than on two (at $1024^3$, 2.08\,s
for the forward and 2.16\,s for the backprojection), while ASTRA keeps gaining (1.68 and 0.96\,s). TIGRE 3.3.0 took
3.37, 2.38 and 2.21\,s for the same forward projection on 1, 2 and 4 GPUs: 3.3.1 is faster on every number of
GPUs, but most on one, because its kernels are faster and the rest of a multi-GPU call is not.

That rest is a part of each multi-GPU call that does not shrink with the number of GPUs
(Table~\ref{tab:lin-fixed}): about 2\,s at $1024^3$ and 0.3--0.4\,s at $512^3$, almost the same for every TIGRE
operator. We could not trace it on this server (it has no Nsight Systems); two parts of it are known. Every GPU
needs a full set of the data it reads: a forward projection splits the angles between the GPUs, and every GPU
gets the whole image (at $1024^3$, 4\,GB to each of them); a backprojection splits the image, and every GPU gets
all the projections. And the caller's arrays are page-locked, so that the copies to the GPUs overlap. On two GPUs
page-locking costs more than it saves, and TIGRE 3.3.1 does not do it there (\#917: 5--21\% faster on two GPUs
than with page-locking, 4--15\% slower on four). The interpolated projector, whose kernels are long, is still 2.5
times faster on four GPUs than on one.

\paragraph{One machine.} All the multi-GPU results in this note come from one server: four Quadro RTX 6000
(Turing, 2018) connected only through PCIe and the host (no NVLink;
\texttt{nvidia-smi topo} reports \texttt{SYS} between every pair). A multi-GPU call depends on how fast the host
feeds the GPUs, more than a one-GPU call does, so on other machines the balance between TIGRE and ASTRA, the
number of GPUs from which TIGRE stops gaining, and the page-locking threshold of \#917 may differ. We would
welcome measurements on other multi-GPU machines, and on other GPUs in general: the suite that produced every
table here is in this note's folder (Code availability), and results can be shared in an issue on TIGRE's GitHub.
\IfFileExists{tables/lin_size_multi.tex}{\begin{table}[H]
\centering
\caption{Linux, Quadro RTX 6000. Cone beam size sweep on 1, 2 and 4 GPUs (NumPy arrays). Speed-up: time on 1 GPU over time on $g$ GPUs, for each toolbox on its own.}
\label{tab:lin-size_multi}
\small\setlength{\tabcolsep}{3.5pt}
\begin{adjustbox}{max width=\linewidth}
\begin{tabular}{llrrrrrrr}
\toprule
& & \multicolumn{3}{c}{TIGRE/ASTRA ratio} & \multicolumn{2}{c}{TIGRE speed-up} & \multicolumn{2}{c}{ASTRA speed-up}\\
\cmidrule(lr){3-5}\cmidrule(lr){6-7}\cmidrule(lr){8-9}
problem & operation & 1 GPU & 2 GPUs & 4 GPUs & 2 GPUs & 4 GPUs & 2 GPUs & 4 GPUs\\
\midrule
128$^3$ & fp, Siddon & \tfast{$0.64\pm0.09$} & \tfast{$0.90\pm0.09$} & \tslow{$1.40\pm0.05$} & 0.78 & 0.43 & 1.12 & 0.96\\
 & fp, interpolated & \tslow{$2.06\pm0.03$} & \tslow{$1.62\pm0.45$} & \tslow{$1.66\pm0.11$} & 1.49 & 1.21 & 1.12 & 0.96\\
 & bp, matched & \tfast{$0.72\pm0.01$} & \tslow{$1.39\pm0.04$} & \tslow{$2.18\pm0.06$} & 0.64 & 0.30 & 1.24 & 0.90\\
 & bp, FDK weights & \tfast{$0.71\pm0.01$} & \tslow{$1.39\pm0.04$} & \tslow{$2.17\pm0.11$} & 0.63 & 0.30 & 1.24 & 0.90\\
\addlinespace
256$^3$ & fp, Siddon & \tfast{$0.75\pm0.02$} & \tslow{$1.13\pm0.08$} & \tslow{$1.68\pm0.10$} & 0.87 & 0.68 & 1.33 & 1.50\\
 & fp, interpolated & \tslow{$2.29\pm0.01$} & \tslow{$1.79\pm0.12$} & \tslow{$1.79\pm0.21$} & 1.69 & 1.99 & 1.33 & 1.50\\
 & bp, matched & \tfast{$0.99\pm0.00$} & \tslow{$1.46\pm0.06$} & \tslow{$2.49\pm0.04$} & 1.01 & 0.73 & 1.50 & 1.83\\
 & bp, FDK weights & \tfast{$0.89\pm0.01$} & \tslow{$1.45\pm0.06$} & \tslow{$2.48\pm0.04$} & 0.91 & 0.66 & 1.50 & 1.83\\
\addlinespace
512$^3$ & fp, Siddon & \tfast{$0.46\pm0.00$} & \tfast{$0.76\pm0.03$} & \tslow{$1.53\pm0.05$} & 0.88 & 0.60 & 1.45 & 1.98\\
 & fp, interpolated & \tslow{$2.60\pm0.02$} & \tslow{$2.15\pm0.04$} & \tslow{$2.22\pm0.14$} & 1.76 & 2.31 & 1.45 & 1.98\\
 & bp, matched & \tslow{$1.09\pm0.01$} & \tslow{$1.50\pm0.26$} & \tslow{$2.50\pm0.14$} & 1.26 & 1.10 & 1.69 & 2.53\\
 & bp, FDK weights & \tfast{$0.97\pm0.01$} & \tslow{$1.44\pm0.25$} & \tslow{$2.60\pm0.15$} & 1.17 & 0.95 & 1.69 & 2.53\\
\addlinespace
1024$^3$ & fp, Siddon & \tfast{$0.48\pm0.00$} & \tfast{$0.66\pm0.01$} & \tslow{$1.26\pm0.12$} & 1.24 & 0.86 & 1.73 & 2.24\\
 & fp, interpolated & \tslow{$2.68\pm0.02$} & \tslow{$2.65\pm0.05$} & \tslow{$2.43\pm0.21$} & 1.75 & 2.50 & 1.73 & 2.24\\
 & bp, matched & \tslow{$1.14\pm0.01$} & \tslow{$1.37\pm0.02$} & \tslow{$2.25\pm0.06$} & 1.28 & 1.41 & 1.55 & 2.78\\
 & bp, FDK weights & \tfast{$0.99\pm0.00$} & \tslow{$1.28\pm0.02$} & \tslow{$2.15\pm0.05$} & 1.20 & 1.28 & 1.55 & 2.78\\
\bottomrule
\end{tabular}
\end{adjustbox}
\end{table}
}{}
\IfFileExists{tables/lin_parallel_multi.tex}{\begin{table}[H]
\centering
\caption{Linux, Quadro RTX 6000. Parallel beam on 1, 2 and 4 GPUs (NumPy arrays).}
\label{tab:lin-parallel_multi}
\small\setlength{\tabcolsep}{3.5pt}
\begin{adjustbox}{max width=\linewidth}
\begin{tabular}{llrrrrrrr}
\toprule
& & \multicolumn{3}{c}{TIGRE/ASTRA ratio} & \multicolumn{2}{c}{TIGRE speed-up} & \multicolumn{2}{c}{ASTRA speed-up}\\
\cmidrule(lr){3-5}\cmidrule(lr){6-7}\cmidrule(lr){8-9}
problem & operation & 1 GPU & 2 GPUs & 4 GPUs & 2 GPUs & 4 GPUs & 2 GPUs & 4 GPUs\\
\midrule
256$^3$ & fp, Siddon & \tfast{$0.95\pm0.02$} & \tslow{$1.44\pm0.04$} & \tslow{$2.05\pm0.09$} & 0.84 & 0.63 & 1.29 & 1.36\\
 & fp, interpolated & \tslow{$2.66\pm0.01$} & \tslow{$2.04\pm0.06$} & \tslow{$2.12\pm0.05$} & 1.68 & 1.72 & 1.29 & 1.36\\
 & bp & \tfast{$0.87\pm0.00$} & \tslow{$1.53\pm0.06$} & \tslow{$2.65\pm0.07$} & 0.88 & 0.63 & 1.55 & 1.92\\
\addlinespace
512$^3$ & fp, Siddon & \tfast{$0.60\pm0.00$} & \tslow{$1.00\pm0.04$}$^\dagger$ & \tslow{$1.84\pm0.07$} & 0.86 & 0.59 & 1.46 & 1.80\\
 & fp, interpolated & \tslow{$2.96\pm0.03$} & \tslow{$2.45\pm0.02$} & \tslow{$2.39\pm0.18$} & 1.76 & 2.25 & 1.46 & 1.80\\
 & bp & \tfast{$0.98\pm0.01$} & \tslow{$1.67\pm0.10$} & \tslow{$2.79\pm0.10$} & 1.11 & 0.90 & 1.91 & 2.59\\
\addlinespace
1024$^3$ & fp, Siddon & \tfast{$0.65\pm0.00$} & \tfast{$0.86\pm0.02$} & \tslow{$1.56\pm0.13$} & 1.23 & 0.83 & 1.61 & 1.98\\
 & fp, interpolated & \tslow{$3.26\pm0.04$} & \tslow{$3.04\pm0.05$} & \tslow{$2.70\pm0.22$} & 1.73 & 2.39 & 1.61 & 1.98\\
 & bp & \tfast{$0.98\pm0.02$} & \tslow{$1.28\pm0.01$} & \tslow{$2.30\pm0.01$} & 1.17 & 1.20 & 1.54 & 2.81\\
\bottomrule
\end{tabular}
\end{adjustbox}
\end{table}
}{}
\IfFileExists{tables/lin_fixed.tex}{\begin{table}[H]
\centering
\caption{Linux, Quadro RTX 6000. Each call's time split into a part that does not shrink with more GPUs (``fixed'') and a part that does, fitted as $t_g = \text{fixed} + \text{parallel}/g$ to the 1- and 4-GPU times (NumPy arrays). The measured 2-GPU time $t_2$ checks the fit.}
\label{tab:lin-fixed}
\small\setlength{\tabcolsep}{4pt}
\begin{adjustbox}{max width=\linewidth}
\begin{tabular}{llrrrrrr}
\toprule
problem & call & $t_1$ ms & $t_4$ ms & fixed ms & fixed \% & $t_2$ ms & $t_2$ fit ms\\
\midrule
cone 512$^3$ & TIGRE fp Siddon & 218 & 366 & 415 & 191 & 249 & 316\\
 & TIGRE fp interpolated & 1238 & 536 & 302 & 24 & 704 & 770\\
 & TIGRE bp matched & 373 & 338 & 326 & 87 & 297 & 350\\
 & ASTRA fp & 476 & 240 & 162 & 34 & 328 & 319\\
 & ASTRA bp & 344 & 136 & 67 & 19 & 204 & 206\\
\addlinespace
cone 1024$^3$ & TIGRE fp Siddon & 1790 & 2075 & 2170 & 121 & 1441 & 1980\\
 & TIGRE fp interpolated & 10080 & 4039 & 2025 & 20 & 5772 & 6053\\
 & TIGRE bp matched & 3030 & 2155 & 1864 & 61 & 2358 & 2447\\
 & TIGRE bp FDK & 2638 & 2056 & 1861 & 71 & 2196 & 2250\\
 & ASTRA fp & 3754 & 1679 & 988 & 26 & 2170 & 2371\\
 & ASTRA bp & 2664 & 959 & 391 & 15 & 1719 & 1527\\
\addlinespace
parallel 1024$^3$ & TIGRE fp Siddon & 1739 & 2099 & 2218 & 128 & 1418 & 1979\\
 & TIGRE fp interpolated & 8685 & 3628 & 1942 & 22 & 5016 & 5314\\
 & TIGRE bp & 2464 & 2053 & 1916 & 78 & 2100 & 2190\\
 & ASTRA fp & 2655 & 1339 & 900 & 34 & 1645 & 1778\\
 & ASTRA bp & 2512 & 893 & 353 & 14 & 1636 & 1432\\
\bottomrule
\end{tabular}
\end{adjustbox}
\end{table}
}{}

\subsection{Large problems: $N^3$ images, $N$ angles}\label{sec:large}
TIGRE's multi-GPU code was designed for a different regime~\cite{tigre2020multigpu}: as many angles as voxels
along a side, which gives several times more kernel work for each byte moved, and images larger than one GPU,
which TIGRE splits between the GPUs and through host memory. Table~\ref{tab:paper} measures it on the Linux
server, with NumPy arrays. At $1536^3$ the image takes 13.6\,GB and the projections as much, more than one GPU's
24\,GB together, and both toolboxes split the problem. TIGRE's Siddon forward projection is 4.1 times faster than
ASTRA's on one GPU and 2.3 times on four, and its FDK-weighted backprojection 1.4 times faster on one GPU. Extra
GPUs help ASTRA more: from one GPU to four its backprojection becomes 3.4 times faster and TIGRE's 2.0 times, so
on four GPUs ASTRA's backprojection is faster (8.3\,s, against 12.5\,s for TIGRE's matched one).

\begin{table}[H]
\centering\small
\caption{Linux, Quadro RTX 6000. Cone beam, $N^3$ image, $N^2$ detector covering the image, $N$ angles, NumPy
arrays: time of one call in s (median of three rounds). At $1536^3$ the image and the projections do not fit in
one GPU together, and both toolboxes split them.}
\label{tab:paper}
\begin{adjustbox}{max width=\linewidth}
\begin{tabular}{rrrrrrr}
\toprule
$N$ & GPUs & TIGRE fp Siddon & ASTRA fp & TIGRE bp matched & TIGRE bp FDK & ASTRA bp \\
\midrule
512 & 1 & 0.29 & 0.66 & 0.44 & 0.39 & 0.45 \\
 & 2 & 0.31 & 0.40 & 0.33 & 0.31 & 0.26 \\
 & 4 & 0.43 & 0.29 & 0.39 & 0.41 & 0.17 \\
\addlinespace[2pt]
1024 & 1 & 4.15 & 9.82 & 5.56 & 4.50 & 5.96 \\
 & 2 & 2.67 & 5.58 & 3.62 & 3.14 & 3.37 \\
 & 4 & 3.14 & 3.30 & 3.23 & 2.95 & 1.78 \\
\addlinespace[2pt]
1536 & 1 & 20.50 & 83.82 & 25.57 & 20.14 & 28.51 \\
 & 2 & 12.53 & 43.09 & 15.74 & 13.15 & 15.51 \\
 & 4 & 10.35 & 24.16 & 12.53 & 11.10 & 8.34 \\
\bottomrule
\end{tabular}
\end{adjustbox}
\end{table}

\section{Where ASTRA is faster, and why}\label{sec:why}
On one GPU, every one of ASTRA's wins falls in one of four causes. On the RTX 4070 they are the interpolated
projector (51 comparisons), the 2D Siddon projector (9) and the backprojection of the smallest images with GPU
arrays (4). On the Quadro, the same three (52, 10 and 5) and the backprojection of large images from host
arrays (14), whose cause we have narrowed down but not confirmed. On several GPUs, ASTRA gains more from extra
GPUs, for the reasons given in the previous section.

\paragraph{The interpolated forward projector (1.0--8.3$\times$).}
TIGRE's interpolated projector samples the image with trilinear 3D texture fetches every
\texttt{geo.accuracy}\,=\,0.5 voxels along the ray. ASTRA's projector (Joseph's method) takes one sample per
slice crossed, with a bilinear fetch from a layered 2D texture. A trilinear fetch costs twice a bilinear one on
NVIDIA GPUs (the effect that made TIGRE's backprojectors faster in 3.3.1, Section~\ref{sec:changes}), and at the
default step TIGRE takes about twice as many samples. Table~\ref{tab:interp} shows the cost of the step: at
one sample per voxel TIGRE is still 1.7--4 times slower. A Joseph-type kernel would close the gap; it was not
part of this release. TIGRE's Siddon projector, which computes exact intersection lengths, is the faster
choice in 3D.
\begin{table}[H]
\centering\small
\caption{TIGRE's interpolated forward projection against its sampling step (\texttt{geo.accuracy}, in voxels),
fan beam, PyTorch tensors, ms.}
\label{tab:interp}
\begin{tabular}{lrrrrrr}
\toprule
problem & 0.25 & 0.5 (default) & 1 & 2 & TIGRE Siddon & ASTRA \\
\midrule
$512^2$, 360 angles & 8.84 & 4.71 & 2.67 & 1.43 & 1.45 & 1.55 \\
$1024^2$, 720 angles & 55.9 & 28.8 & 15.2 & 5.07 & 4.81 & 3.69 \\
\bottomrule
\end{tabular}
\end{table}

\paragraph{The Siddon forward projector in 2D (1.0--1.3$\times$ on the RTX 4070, 1.5--2.4$\times$ on the Quadro).}
In 3D, TIGRE's Siddon projector takes 0.3--0.9 of ASTRA's time; in 2D fan beam it is slower. The kernel is: in
Nsight Systems, at $1024^2$ with 720 angles, TIGRE's kernels take 2.9\,ms and ASTRA's 1.7\,ms on the RTX 4070.
Siddon's method visits every voxel boundary the ray crosses, about $1.3N$ steps for an $N\times N$ image against
the $N$ samples of ASTRA's 2D projector (Joseph's method), and each step chooses the next boundary. In 3D the
same comparison favours TIGRE, because ASTRA's 3D projector pays for it in other ways (Section 5.5 of the Siddon
note~\cite{tigre2025siddon}). Reading the image through a layered 2D texture, the change that halved the
backprojectors' cost, made the Siddon projector 1.3--1.5 times \emph{slower}, and was not kept. The Quadro, an
older GPU, is further behind ASTRA here.

\paragraph{Backprojection of very small images with GPU arrays ($32^3$ and $64^3$).}
With GPU tensors, TIGRE's backprojection of a $32^3$ image takes 0.94\,ms and ASTRA's 0.55\,ms on the RTX 4070.
In an earlier run, with 0.93 and 0.62\,ms, Nsight Systems
shows 0.09\,ms of kernels in both; TIGRE spends about 0.23\,ms in Python before the call (checking and copying
the geometry, kept on purpose: it catches wrong geometries before they reach the GPU) and about 0.2\,ms more in
CUDA calls than ASTRA (the projections' copy into the texture array and the per-kernel constants). From $128^3$
TIGRE is faster.

\paragraph{Backprojection from host arrays on the Quadro (1.02--1.75$\times$).}
On the Linux server, TIGRE's backprojection with GPU arrays takes 0.6--0.9 of ASTRA's time, but with NumPy arrays
it is slower from $512^3$, most with few angles ($512^3$, 45 angles: 240\,ms against 137\,ms). The difference is
in what a call with host arrays adds to one with GPU arrays (Table~\ref{tab:hostcost}): for TIGRE about 207\,ms at
$512^3$ whatever the number of angles, for ASTRA 106\,ms with 45 angles, growing with the projections it copies.
On the RTX 4070 both add about 80\,ms. TIGRE's added time does not depend on the projections, so it is in the
result: allocating the $512^3$ image on the host (512\,MB) and copying it back from the GPU. A likely cause is the
host memory itself. TIGRE allocates the result with \texttt{malloc}, so the copy into it touches 131072 fresh
4\,kB pages; ASTRA writes into an array NumPy allocated, which on Linux asks for huge pages. That would also
explain why Windows does not show it. We have not yet confirmed it.

\begin{table}[H]
\centering\small
\caption{What a call with NumPy arrays adds to the same call with GPU tensors: $t_\text{host}-t_\text{tensors}$
in ms, cone $512^3$, one GPU.}
\label{tab:hostcost}
\begin{tabular}{llrrr}
\toprule
GPU & angles & TIGRE bp matched & TIGRE bp FDK & ASTRA bp \\
\midrule
Quadro RTX 6000 & 45 & 212 & 218 & 106 \\
 & 180 & 206 & 209 & 115 \\
 & 360 & 207 & 211 & 128 \\
 & 1440 & 207 & 207 & 218 \\
\addlinespace[2pt]
RTX 4070 & 45 & 79 & 80 & 81 \\
 & 180 & 92 & 77 & 99 \\
 & 360 & 77 & 83 & 117 \\
 & 1440 & 76 & 98 & 273 \\
\bottomrule
\end{tabular}
\end{table}

\section{Discussion}
On one GPU, TIGRE 3.3.1's Siddon projector is faster than ASTRA's for every 3D problem, and both backprojectors
for every 2D problem and every 3D one from $128^3$ ($256^3$ on the Quadro). The exception is the Linux server,
where the backprojectors are slower than ASTRA's with host arrays from $512^3$ (Section~\ref{sec:why}). The changes that
brought this were of two kinds. Most of 3.3.1's work removed fixed costs that did not grow with the problem and
that had dominated small and 2D problems. The largest single change was not in the host code: reading the
projections through a layered 2D texture instead of a 3D one halves the cost of every texture fetch in the
backprojectors, without changing a bit of the results. ASTRA's remaining advantages come from an algorithmic
choice (Joseph's interpolation against Siddon's exact lengths, in 2D), from the fixed cost of a call on the
smallest images, or from the interpolated projector, which this release left as it was.

On several GPUs, TIGRE gains less than ASTRA from each extra GPU for problems that fit in one. Two GPUs are now
faster than one for large images, but four are slower than two, because every multi-GPU call copies a full set of
its input to every GPU. That part is the next to reduce, without giving up what TIGRE's multi-GPU path is for:
splitting the image and the projections across GPUs and through host memory, so that it reconstructs images
larger than one GPU, on one GPU or on several~\cite{tigre2020multigpu}. In that regime TIGRE's forward projection
is 2.3--4.1 times faster than ASTRA's, its backprojection faster on one GPU and slower on four
(Section~\ref{sec:large}). For a problem that fits in one GPU, one or two GPUs (\texttt{gpuids}) are the fastest
choice in TIGRE 3.3.1 on the server we measured; with one multi-GPU machine, these conclusions need checking on
others (Section~\ref{sec:results}, ``One machine'').

\appendix
\section{Per-case results on the Linux server}\label{app:linux}
One GPU of the Linux server (Quadro RTX 6000, Section~\ref{sec:setup}); same suite and layout as the tables of
Section~\ref{sec:results}.
\IfFileExists{tables/lin_size1.tex}{\begin{table}[H]
\centering
\caption{Linux, Quadro RTX 6000. Cone beam, $n^3$ image, $n^2$ detector covering the image (409.6\,mm), 360 angles, 1 GPU.}
\label{tab:lin-size1}
\small\setlength{\tabcolsep}{4pt}
\begin{adjustbox}{max width=\linewidth}
\begin{tabular}{llrrrrrr}
\toprule
& & \multicolumn{3}{c}{NumPy arrays, 1 GPU} & \multicolumn{3}{c}{PyTorch GPU tensors, 1 GPU}\\
\cmidrule(lr){3-5}\cmidrule(lr){6-8}
problem & operation & TIGRE ms & ASTRA ms & ratio & TIGRE ms & ASTRA ms & ratio\\
\midrule
128$^3$ & fp, Siddon & 6.72 & 10.8 & \tfast{$0.64\pm0.09$} & 4.96 & 6.73 & \tfast{$0.78\pm0.18$}\\
 & fp, interpolated & 22.3 & 10.8 & \tslow{$2.06\pm0.03$} & 21.4 & 6.73 & \tslow{$3.17\pm0.05$}\\
 & bp, matched & 6.16 & 8.56 & \tfast{$0.72\pm0.01$} & 4.31 & 4.75 & \tfast{$0.91\pm0.00$}\\
 & bp, FDK weights & 6.07 & 8.56 & \tfast{$0.71\pm0.01$} & 3.74 & 4.75 & \tfast{$0.79\pm0.00$}\\
\addlinespace
256$^3$ & fp, Siddon & 51.4 & 69.0 & \tfast{$0.75\pm0.02$} & 23.6 & 44.2 & \tfast{$0.53\pm0.00$}\\
 & fp, interpolated & 158 & 69.0 & \tslow{$2.29\pm0.01$} & 154 & 44.2 & \tslow{$3.48\pm0.02$}\\
 & bp, matched & 48.9 & 49.6 & \tfast{$0.99\pm0.00$} & 22.5 & 28.0 & \tfast{$0.80\pm0.01$}\\
 & bp, FDK weights & 44.0 & 49.6 & \tfast{$0.89\pm0.01$} & 17.6 & 28.0 & \tfast{$0.63\pm0.01$}\\
\addlinespace
512$^3$ & fp, Siddon & 218 & 476 & \tfast{$0.46\pm0.00$} & 176 & 361 & \tfast{$0.49\pm0.01$}\\
 & fp, interpolated & 1238 & 476 & \tslow{$2.60\pm0.02$} & 1204 & 361 & \tslow{$3.33\pm0.02$}\\
 & bp, matched & 373 & 344 & \tslow{$1.09\pm0.01$} & 172 & 216 & \tfast{$0.79\pm0.01$}\\
 & bp, FDK weights & 335 & 344 & \tfast{$0.97\pm0.01$} & 130 & 216 & \tfast{$0.60\pm0.00$}\\
\addlinespace
1024$^3$ & fp, Siddon & 1790 & 3754 & \tfast{$0.48\pm0.00$} & 1469 & 3099 & \tfast{$0.47\pm0.01$}\\
 & fp, interpolated & 10080 & 3754 & \tslow{$2.68\pm0.02$} & 9768 & 3099 & \tslow{$3.14\pm0.05$}\\
 & bp, matched & 3030 & 2664 & \tslow{$1.14\pm0.01$} & 1426 & 1722 & \tfast{$0.83\pm0.00$}\\
 & bp, FDK weights & 2638 & 2664 & \tfast{$0.99\pm0.00$} & 1024 & 1722 & \tfast{$0.59\pm0.00$}\\
\bottomrule
\end{tabular}
\end{adjustbox}
\end{table}
}{}
\IfFileExists{tables/lin_parallel1.tex}{\begin{table}[H]
\centering
\caption{Linux, Quadro RTX 6000. Parallel beam (TIGRE parallel mode, ASTRA \texttt{parallel3d}), $n^3$, 360 angles, 1 GPU.}
\label{tab:lin-parallel1}
\small\setlength{\tabcolsep}{4pt}
\begin{adjustbox}{max width=\linewidth}
\begin{tabular}{llrrrrrr}
\toprule
& & \multicolumn{3}{c}{NumPy arrays, 1 GPU} & \multicolumn{3}{c}{PyTorch GPU tensors, 1 GPU}\\
\cmidrule(lr){3-5}\cmidrule(lr){6-8}
problem & operation & TIGRE ms & ASTRA ms & ratio & TIGRE ms & ASTRA ms & ratio\\
\midrule
256$^3$ & fp, Siddon & 51.3 & 54.5 & \tfast{$0.95\pm0.02$} & 24.3 & 28.9 & \tfast{$0.84\pm0.01$}\\
 & fp, interpolated & 145 & 54.5 & \tslow{$2.66\pm0.01$} & 136 & 28.9 & \tslow{$4.72\pm0.01$}\\
 & bp & 42.6 & 49.2 & \tfast{$0.87\pm0.00$} & 15.2 & 26.6 & \tfast{$0.57\pm0.00$}\\
\addlinespace
512$^3$ & fp, Siddon & 218 & 366 & \tfast{$0.60\pm0.00$} & 173 & 252 & \tfast{$0.68\pm0.01$}\\
 & fp, interpolated & 1085 & 366 & \tslow{$2.96\pm0.03$} & 1041 & 252 & \tslow{$4.12\pm0.02$}\\
 & bp & 319 & 326 & \tfast{$0.98\pm0.01$} & 107 & 196 & \tfast{$0.54\pm0.00$}\\
\addlinespace
1024$^3$ & fp, Siddon & 1739 & 2655 & \tfast{$0.65\pm0.00$} & 1411 & 2004 & \tfast{$0.70\pm0.01$}\\
 & fp, interpolated & 8685 & 2655 & \tslow{$3.26\pm0.04$} & 8361 & 2004 & \tslow{$4.16\pm0.06$}\\
 & bp & 2464 & 2512 & \tfast{$0.98\pm0.02$} & 838 & 1566 & \tfast{$0.54\pm0.00$}\\
\bottomrule
\end{tabular}
\end{adjustbox}
\end{table}
}{}
\IfFileExists{tables/lin_fan.tex}{\begin{table}[H]
\centering
\caption{Linux, Quadro RTX 6000. Fan beam (2D): TIGRE's one-slice cone geometry against ASTRA's \texttt{fanflat} (FP\_CUDA, BP\_CUDA, FBP\_CUDA), $n^2$ image, $n$ detector pixels, 1 GPU.}
\label{tab:lin-fan}
\small\setlength{\tabcolsep}{3pt}
\begin{adjustbox}{max width=\linewidth}
\begin{tabular}{llrrrrrr}
\toprule
& & \multicolumn{3}{c}{NumPy arrays, 1 GPU} & \multicolumn{3}{c}{PyTorch GPU tensors, 1 GPU}\\
\cmidrule(lr){3-5}\cmidrule(lr){6-8}
problem & operation & TIGRE ms & ASTRA ms & ratio & TIGRE ms & ASTRA ms & ratio\\
\midrule
512$^2$ & fp, Siddon & 2.50 & 1.11 & \tslow{$2.39\pm0.60$} & 2.22 & 1.01 & \tslow{$2.20\pm0.00$}\\
 & fp, interpolated & 4.64 & 1.11 & \tslow{$4.19\pm0.04$} & 4.36 & 1.01 & \tslow{$4.31\pm0.02$}\\
 & bp, matched & 2.48 & 3.33 & \tfast{$0.74\pm0.00$} & 1.99 & 3.18 & \tfast{$0.63\pm0.00$}\\
 & bp, FDK weights & 2.45 & 3.33 & \tfast{$0.73\pm0.00$} & 1.96 & 3.18 & \tfast{$0.62\pm0.00$}\\
\addlinespace
1024$^2$ & fp, Siddon & 3.82 & 2.28 & \tslow{$1.68\pm0.01$} & 3.11 & 1.95 & \tslow{$1.60\pm0.01$}\\
 & fp, interpolated & 11.6 & 2.28 & \tslow{$5.10\pm0.02$} & 10.9 & 1.95 & \tslow{$5.59\pm0.01$}\\
 & bp, matched & 3.99 & 10.7 & \tfast{$0.37\pm0.00$} & 3.16 & 10.3 & \tfast{$0.31\pm0.00$}\\
 & bp, FDK weights & 3.87 & 10.7 & \tfast{$0.36\pm0.00$} & 3.03 & 10.3 & \tfast{$0.29\pm0.00$}\\
\addlinespace
2048$^2$ & fp, Siddon & 9.44 & 6.40 & \tslow{$1.48\pm0.00$} & 7.33 & 4.93 & \tslow{$1.49\pm0.00$}\\
 & fp, interpolated & 38.6 & 6.40 & \tslow{$6.04\pm0.05$} & 36.5 & 4.93 & \tslow{$7.42\pm0.04$}\\
 & bp, matched & 10.3 & 40.2 & \tfast{$0.26\pm0.00$} & 8.22 & 38.7 & \tfast{$0.21\pm0.00$}\\
 & bp, FDK weights & 9.67 & 40.2 & \tfast{$0.24\pm0.00$} & 7.53 & 38.7 & \tfast{$0.19\pm0.00$}\\
\addlinespace
4096$^2$ & fp, Siddon & 38.9 & 21.7 & \tslow{$1.80\pm0.02$} & 32.2 & 17.3 & \tslow{$1.86\pm0.00$}\\
 & fp, interpolated & 151 & 21.7 & \tslow{$6.98\pm0.04$} & 144 & 17.3 & \tslow{$8.30\pm0.01$}\\
 & bp, matched & 53.9 & 165 & \tfast{$0.33\pm0.00$} & 28.9 & 155 & \tfast{$0.19\pm0.00$}\\
 & bp, FDK weights & 51.8 & 165 & \tfast{$0.31\pm0.00$} & 26.2 & 155 & \tfast{$0.17\pm0.00$}\\
\addlinespace
1024$^2$, 1440 angles & fp, Siddon & 7.79 & 4.86 & \tslow{$1.61\pm0.00$} & 7.16 & 4.56 & \tslow{$1.57\pm0.01$}\\
 & fp, interpolated & 38.1 & 4.86 & \tslow{$7.85\pm0.03$} & 37.5 & 4.56 & \tslow{$8.23\pm0.00$}\\
 & bp, matched & 9.99 & 39.8 & \tfast{$0.25\pm0.00$} & 8.55 & 39.1 & \tfast{$0.22\pm0.00$}\\
 & bp, FDK weights & 9.41 & 39.8 & \tfast{$0.24\pm0.00$} & 7.92 & 39.1 & \tfast{$0.20\pm0.00$}\\
\bottomrule
\end{tabular}
\end{adjustbox}
\end{table}
}{}
\IfFileExists{tables/lin_angles.tex}{\begin{table}[H]
\centering
\caption{Linux, Quadro RTX 6000. Cone beam $512^3$, varying the number of angles, 1 GPU.}
\label{tab:lin-angles}
\small\setlength{\tabcolsep}{4pt}
\begin{adjustbox}{max width=\linewidth}
\begin{tabular}{llrrrrrr}
\toprule
& & \multicolumn{3}{c}{NumPy arrays, 1 GPU} & \multicolumn{3}{c}{PyTorch GPU tensors, 1 GPU}\\
\cmidrule(lr){3-5}\cmidrule(lr){6-8}
problem & operation & TIGRE ms & ASTRA ms & ratio & TIGRE ms & ASTRA ms & ratio\\
\midrule
45 angles & fp, Siddon & 84.0 & 119 & \tfast{$0.71\pm0.01$} & 35.8 & 58.9 & \tfast{$0.61\pm0.00$}\\
 & fp, interpolated & 212 & 119 & \tslow{$1.78\pm0.02$} & 164 & 58.9 & \tslow{$2.78\pm0.01$}\\
 & bp, matched & 240 & 137 & \tslow{$1.75\pm0.00$} & 28.4 & 30.9 & \tfast{$0.92\pm0.02$}\\
 & bp, FDK weights & 237 & 137 & \tslow{$1.73\pm0.01$} & 19.3 & 30.9 & \tfast{$0.62\pm0.01$}\\
\addlinespace
90 angles & fp, Siddon & 98.5 & 165 & \tfast{$0.59\pm0.01$} & 53.2 & 101 & \tfast{$0.53\pm0.00$}\\
 & fp, interpolated & 359 & 165 & \tslow{$2.16\pm0.04$} & 314 & 101 & \tslow{$3.12\pm0.03$}\\
 & bp, matched & 259 & 163 & \tslow{$1.59\pm0.00$} & 48.8 & 56.5 & \tfast{$0.86\pm0.01$}\\
 & bp, FDK weights & 248 & 163 & \tslow{$1.53\pm0.04$} & 35.2 & 56.5 & \tfast{$0.62\pm0.00$}\\
\addlinespace
180 angles & fp, Siddon & 136 & 265 & \tfast{$0.51\pm0.01$} & 94.1 & 187 & \tfast{$0.50\pm0.00$}\\
 & fp, interpolated & 655 & 265 & \tslow{$2.46\pm0.03$} & 613 & 187 & \tslow{$3.28\pm0.02$}\\
 & bp, matched & 297 & 225 & \tslow{$1.32\pm0.00$} & 90.5 & 110 & \tfast{$0.83\pm0.00$}\\
 & bp, FDK weights & 275 & 225 & \tslow{$1.23\pm0.01$} & 66.7 & 110 & \tfast{$0.61\pm0.00$}\\
\addlinespace
360 angles & fp, Siddon & 220 & 477 & \tfast{$0.46\pm0.00$} & 176 & 362 & \tfast{$0.49\pm0.01$}\\
 & fp, interpolated & 1251 & 477 & \tslow{$2.62\pm0.02$} & 1209 & 362 & \tslow{$3.34\pm0.01$}\\
 & bp, matched & 380 & 346 & \tslow{$1.10\pm0.01$} & 173 & 218 & \tfast{$0.79\pm0.01$}\\
 & bp, FDK weights & 342 & 346 & \tfast{$0.99\pm0.02$}$^\dagger$ & 131 & 218 & \tfast{$0.60\pm0.01$}\\
\addlinespace
720 angles & fp, Siddon & 387 & 933 & \tfast{$0.42\pm0.01$} & 343 & 732 & \tfast{$0.47\pm0.01$}\\
 & fp, interpolated & 2452 & 933 & \tslow{$2.63\pm0.01$} & 2411 & 732 & \tslow{$3.29\pm0.02$}\\
 & bp, matched & 540 & 607 & \tfast{$0.89\pm0.00$} & 337 & 430 & \tfast{$0.78\pm0.01$}\\
 & bp, FDK weights & 463 & 607 & \tfast{$0.77\pm0.02$} & 258 & 430 & \tfast{$0.60\pm0.00$}\\
\addlinespace
1440 angles & fp, Siddon & 726 & 1822 & \tfast{$0.40\pm0.01$} & 676 & 1495 & \tfast{$0.45\pm0.00$}\\
 & fp, interpolated & 4852 & 1822 & \tslow{$2.66\pm0.01$} & 4813 & 1495 & \tslow{$3.22\pm0.01$}\\
 & bp, matched & 872 & 1075 & \tfast{$0.81\pm0.00$} & 665 & 858 & \tfast{$0.78\pm0.00$}\\
 & bp, FDK weights & 721 & 1075 & \tfast{$0.67\pm0.00$} & 514 & 858 & \tfast{$0.60\pm0.00$}\\
\bottomrule
\end{tabular}
\end{adjustbox}
\end{table}
}{}
\IfFileExists{tables/lin_shape.tex}{\begin{table}[H]
\centering
\caption{Linux, Quadro RTX 6000. Non-cubic problems, 1 GPU: clinical CBCT ($256\times512\times512$ image, $512\times1024$ detector), fine detector ($256^3$ image, $1024^2$ detector), coarse detector ($512^3$ image, $256^2$ detector), 360 angles.}
\label{tab:lin-shape}
\small\setlength{\tabcolsep}{4pt}
\begin{adjustbox}{max width=\linewidth}
\begin{tabular}{llrrrrrr}
\toprule
& & \multicolumn{3}{c}{NumPy arrays, 1 GPU} & \multicolumn{3}{c}{PyTorch GPU tensors, 1 GPU}\\
\cmidrule(lr){3-5}\cmidrule(lr){6-8}
problem & operation & TIGRE ms & ASTRA ms & ratio & TIGRE ms & ASTRA ms & ratio\\
\midrule
clinical CBCT & fp, Siddon & 332 & 813 & \tfast{$0.41\pm0.01$} & 264 & 632 & \tfast{$0.42\pm0.00$}\\
 & fp, interpolated & 2325 & 813 & \tslow{$2.86\pm0.00$} & 2279 & 632 & \tslow{$3.60\pm0.01$}\\
 & bp, matched & 214 & 243 & \tfast{$0.88\pm0.02$} & 104 & 115 & \tfast{$0.91\pm0.01$}\\
 & bp, FDK weights & 184 & 243 & \tfast{$0.75\pm0.03$} & 74.2 & 115 & \tfast{$0.65\pm0.00$}\\
\addlinespace
fine detector & fp, Siddon & 605 & 920 & \tfast{$0.65\pm0.01$} & 276 & 601 & \tfast{$0.46\pm0.01$}\\
 & fp, interpolated & 2210 & 920 & \tslow{$2.40\pm0.02$} & 2183 & 601 & \tslow{$3.63\pm0.02$}\\
 & bp, matched & 153 & 239 & \tfast{$0.64\pm0.01$} & 47.1 & 52.6 & \tfast{$0.90\pm0.00$}\\
 & bp, FDK weights & 164 & 239 & \tfast{$0.70\pm0.05$} & 40.6 & 52.6 & \tfast{$0.77\pm0.02$}\\
\addlinespace
coarse detector & fp, Siddon & 134 & 165 & \tfast{$0.81\pm0.01$} & 90.3 & 101 & \tfast{$0.90\pm0.01$}\\
 & fp, interpolated & 426 & 165 & \tslow{$2.58\pm0.04$} & 381 & 101 & \tslow{$3.79\pm0.00$}\\
 & bp, matched & 368 & 315 & \tslow{$1.16\pm0.01$} & 160 & 209 & \tfast{$0.77\pm0.01$}\\
 & bp, FDK weights & 336 & 315 & \tslow{$1.07\pm0.01$} & 125 & 209 & \tfast{$0.60\pm0.01$}\\
\bottomrule
\end{tabular}
\end{adjustbox}
\end{table}
}{}
\IfFileExists{tables/lin_offset.tex}{\begin{table}[H]
\centering
\caption{Linux, Quadro RTX 6000. Cone beam $512^3$ with the detector shifted sideways by 102.4\,mm (a quarter of its width; ASTRA \texttt{cone\_vec}), 1 GPU.}
\label{tab:lin-offset}
\small\setlength{\tabcolsep}{4pt}
\begin{adjustbox}{max width=\linewidth}
\begin{tabular}{llrrrrrr}
\toprule
& & \multicolumn{3}{c}{NumPy arrays, 1 GPU} & \multicolumn{3}{c}{PyTorch GPU tensors, 1 GPU}\\
\cmidrule(lr){3-5}\cmidrule(lr){6-8}
problem & operation & TIGRE ms & ASTRA ms & ratio & TIGRE ms & ASTRA ms & ratio\\
\midrule
u offset 102.4 mm & fp, Siddon & 205 & 461 & \tfast{$0.44\pm0.01$} & 145 & 347 & \tfast{$0.42\pm0.00$}\\
 & fp, interpolated & 1244 & 461 & \tslow{$2.69\pm0.02$} & 1203 & 347 & \tslow{$3.48\pm0.02$}\\
 & bp, matched & 373 & 344 & \tslow{$1.08\pm0.00$} & 168 & 216 & \tfast{$0.78\pm0.01$}\\
 & bp, FDK weights & 337 & 344 & \tfast{$0.98\pm0.01$} & 129 & 216 & \tfast{$0.60\pm0.00$}\\
\bottomrule
\end{tabular}
\end{adjustbox}
\end{table}
}{}
\IfFileExists{tables/lin_calib.tex}{\begin{table}[H]
\centering
\caption{Linux, Quadro RTX 6000. Calibration: the ASTRA convention that matches TIGRE, found on a small problem for each beam type, and how well the projections agree under it (relative $L_2$ difference after a least-squares scale). ``Next best'': the best convention that gives measurably different projections.}
\label{tab:lin-calib}
\small
\begin{adjustbox}{max width=\linewidth}
\begin{tabular}{llrrr}
\toprule
beam & ASTRA convention & agree & scale & next best\\
\midrule
cone & $\theta_A = +1\,\theta_T +90^\circ$ & 5.4e-04 & 0.9998 & 1.6e-01\\
parallel & $\theta_A = +1\,\theta_T +270^\circ$, $u$ flipped & 5.1e-04 & 0.9998 & 4.9e-01\\
fan & $\theta_A = -1\,\theta_T +90^\circ$, $u$ flipped & 5.4e-04 & 0.9999 & 1.2e-01\\
offset & detector shift $+1\,u, +1\,v$ & 5.2e-04 & 0.9998 & 3.9e-01\\
\bottomrule
\end{tabular}
\end{adjustbox}
\end{table}
}{}
\IfFileExists{tables/lin_legacy.tex}{\begin{table}[H]
\centering
\caption{Linux, Quadro RTX 6000. The 3.3.0 note's detector (0.8\,mm pixels whatever $n$), 1 GPU. Compare with Table~\ref{tab:size1} at the same $n$.}
\label{tab:lin-legacy}
\small\setlength{\tabcolsep}{4pt}
\begin{adjustbox}{max width=\linewidth}
\begin{tabular}{llrrrrrr}
\toprule
& & \multicolumn{3}{c}{NumPy arrays, 1 GPU} & \multicolumn{3}{c}{PyTorch GPU tensors, 1 GPU}\\
\cmidrule(lr){3-5}\cmidrule(lr){6-8}
problem & operation & TIGRE ms & ASTRA ms & ratio & TIGRE ms & ASTRA ms & ratio\\
\midrule
256$^3$, 0.8 mm & fp, Siddon & 51.2 & 66.6 & \tfast{$0.77\pm0.02$} & 25.1 & 44.1 & \tfast{$0.57\pm0.00$}\\
 & fp, interpolated & 160 & 66.6 & \tslow{$2.41\pm0.00$} & 154 & 44.1 & \tslow{$3.49\pm0.01$}\\
 & bp, matched & 47.9 & 50.1 & \tfast{$0.96\pm0.00$} & 22.1 & 28.1 & \tfast{$0.79\pm0.00$}\\
 & bp, FDK weights & 44.5 & 50.1 & \tfast{$0.89\pm0.01$} & 18.0 & 28.1 & \tfast{$0.64\pm0.01$}\\
\addlinespace
1024$^3$, 0.8 mm & fp, Siddon & 1679 & 3604 & \tfast{$0.47\pm0.00$} & 1355 & 2950 & \tfast{$0.46\pm0.00$}\\
 & fp, interpolated & 11291 & 3604 & \tslow{$3.12\pm0.03$} & 10990 & 2950 & \tslow{$3.71\pm0.05$}\\
 & bp, matched & 2877 & 2547 & \tslow{$1.12\pm0.01$} & 1283 & 1667 & \tfast{$0.77\pm0.00$}\\
 & bp, FDK weights & 2621 & 2547 & \tslow{$1.02\pm0.01$} & 1002 & 1667 & \tfast{$0.60\pm0.00$}\\
\bottomrule
\end{tabular}
\end{adjustbox}
\end{table}
}{}
\IfFileExists{tables/lin_latency.tex}{\begin{table}[H]
\centering
\caption{Linux, Quadro RTX 6000. Small problems (cone beam, $n^3$ image, $n^2$ detector, 100 angles), 1 GPU.}
\label{tab:lin-latency}
\small\setlength{\tabcolsep}{4pt}
\begin{adjustbox}{max width=\linewidth}
\begin{tabular}{llrrrrrr}
\toprule
& & \multicolumn{3}{c}{NumPy arrays, 1 GPU} & \multicolumn{3}{c}{PyTorch GPU tensors, 1 GPU}\\
\cmidrule(lr){3-5}\cmidrule(lr){6-8}
problem & operation & TIGRE ms & ASTRA ms & ratio & TIGRE ms & ASTRA ms & ratio\\
\midrule
32$^3$ & fp, Siddon & 2.20 & 3.17 & \tfast{$0.69\pm0.03$} & 1.73 & 2.09 & \tfast{$0.82\pm0.06$}\\
 & fp, interpolated & 2.81 & 3.17 & \tfast{$0.86\pm0.06$} & 2.31 & 2.09 & \tslow{$1.10\pm0.07$}\\
 & bp, matched & 2.16 & 2.40 & \tfast{$0.90\pm0.01$} & 1.63 & 1.45 & \tslow{$1.15\pm0.07$}\\
 & bp, FDK weights & 2.15 & 2.40 & \tfast{$0.89\pm0.00$} & 1.62 & 1.45 & \tslow{$1.15\pm0.11$}\\
\addlinespace
64$^3$ & fp, Siddon & 2.26 & 3.75 & \tfast{$0.59\pm0.06$} & 1.72 & 2.62 & \tfast{$0.66\pm0.00$}\\
 & fp, interpolated & 3.00 & 3.75 & \tfast{$0.83\pm0.12$} & 2.64 & 2.62 & \tslow{$1.01\pm0.00$}\\
 & bp, matched & 2.13 & 2.37 & \tfast{$0.90\pm0.00$} & 1.66 & 1.46 & \tslow{$1.13\pm0.00$}\\
 & bp, FDK weights & 2.12 & 2.37 & \tfast{$0.90\pm0.00$} & 1.62 & 1.46 & \tslow{$1.11\pm0.00$}\\
\addlinespace
128$^3$ & fp, Siddon & 3.71 & 6.91 & \tfast{$0.54\pm0.00$} & 2.46 & 4.75 & \tfast{$0.52\pm0.00$}\\
 & fp, interpolated & 8.07 & 6.91 & \tslow{$1.17\pm0.00$} & 7.10 & 4.75 & \tslow{$1.50\pm0.00$}\\
 & bp, matched & 3.64 & 4.18 & \tfast{$0.87\pm0.00$} & 2.41 & 2.26 & \tslow{$1.07\pm0.00$}\\
 & bp, FDK weights & 3.55 & 4.18 & \tfast{$0.85\pm0.00$} & 2.21 & 2.26 & \tfast{$0.98\pm0.00$}\\
\bottomrule
\end{tabular}
\end{adjustbox}
\end{table}
}{}

\section*{Code availability}
TIGRE: \url{https://github.com/CERN/TIGRE}, release 3.3.1. The comparison suite is the file
\path{tigre_vs_astra_suite.py} in \path{docs/technical-notes/speed-3.3.1/}, with the raw results of this note
(\path{results/}) and the scripts that make its tables; its \path{README.md} says how to run it on another machine.
Results from other GPUs and multi-GPU machines are welcome as an issue on TIGRE's GitHub.

\section*{Use of generative AI}
This work and this note were done with substantial help from Claude (Anthropic), a large language model, used
as a coding and writing assistant (Claude Opus, through Claude Code). Under the author's
direction, Claude carried out much of the investigation and profiling, wrote the changes to TIGRE's host code and
their tests, wrote the benchmark suite and the scripts that ran it, and wrote the first draft of this note.

The author conceived and directed the work, decided what to change and how, reviewed and merged every change,
ran the benchmarks on GPU hardware (the results reported here), checked the results, and revised the text. The
author takes full responsibility for the content of this note, including any errors. Claude is not an author: it
cannot take responsibility for the work or for its publication.

\begin{thebibliography}{9}
\bibitem{tigre2016} A.~Biguri, M.~Dosanjh, S.~Hancock, M.~Soleimani, ``TIGRE: a MATLAB-GPU toolbox for CBCT
image reconstruction,'' \emph{Biomedical Physics \& Engineering Express} 2(5):055010, 2016.
doi:10.1088/2057-1976/2/5/055010.
\bibitem{tigre2020multigpu} A.~Biguri, R.~J.~Lindroos, R.~Bryll, H.~Towsyfyan, H.~Deyhle, I.~Harrane,
R.~Boardman, M.~Mavrogordato, M.~Dosanjh, S.~Hancock, T.~Blumensath, ``Arbitrarily large tomography with
iterative algorithms on multiple GPUs using the TIGRE toolbox,'' \emph{Journal of Parallel and Distributed
Computing} 146:52--63, 2020. doi:10.1016/j.jpdc.2020.07.004.
\bibitem{tigre2025v3} A.~Biguri, T.~Sadakane, R.~Lindroos, Y.~Liu, M.~Sabat\'e Landman, Y.~Du, M.~Lohvithee,
S.~Kaser, S.~Hatamikia, R.~Bryll, E.~Valat, S.~Wonglee, T.~Blumensath, C.-B.~Sch\"onlieb, ``TIGRE v3: efficient
and easy to use iterative computed tomographic reconstruction toolbox for real datasets,''
\emph{Engineering Research Express} 7(1):015011, 2025. doi:10.1088/2631-8695/adbb3a.
\bibitem{astra2016} W.~van~Aarle, W.~J.~Palenstijn, J.~Cant, E.~Janssens, F.~Bleichrodt, A.~Dabravolski,
J.~De~Beenhouwer, K.~J.~Batenburg, J.~Sijbers, ``Fast and flexible X-ray tomography using the ASTRA toolbox,''
\emph{Optics Express} 24(22):25129--25147, 2016.
\bibitem{tigre2025siddon} A.~Biguri, ``Siddon--Jacobs ray tracing on GPUs: a technical note on a consistent and
accurate implementation,'' TIGRE technical note, November 2026, revised for TIGRE 3.3.1.
\end{thebibliography}
\end{document}
